% !TEX root = ../main.tex

\section{Introduction}

Fluid flow interaction with coastal bathymetry plays a significant role in mass and momentum transfer, which fundamentally drives the movement of nutrients and tracers throughout the ocean \cite{pawlak03}. Features such as headlands and capes are also capable of developing eddies that can alter the transport of sediment and dissolved pollutants \cite{bastos03}. Flow-bathymetry interactions therefore have important implications the overall circulation regime of the near-coastal ocean system, and understanding its complex physical dynamics may pave light to new mechanisms relating to transport mechanics in biophysical contexts.

In coastal regions with active tidal currents and rough bathymetry, form drag transfers energy and momentum from large scale to submesoscale. Similarly, tides flowing over bathymetry generate internal waves and eddies that transport mass and removes energy from coastal currents \cite{warner14}. There is an extensive collection of studies that observe flow interactions with seamounts, headlands, and capes, which are influenced by form drag generated from assymetric bottom pressure across topography \cite{chor25}. In addition to drag, energy dissipates by wave generation, and wave breaking \cite{klymak21, zemskova21}. 

The eddies which generate as a result of bathymetry-flow interactions are variable both spatially and temporally. Furthermore, vorticity generation is influenced particularly by torque from stress vertical divergence at bottom bathymetry, and bathymetry induces an intensification and redistribution of vorticity when surrounding waters move towards it \cite{jagannathan20}. Flow over bathymetry also generate submesoscale structures, which further complicate dynamics in regards to energy and momentum transfer \cite{srinivasan21}. 

In the context of flow-bathymetry interactions, shoals are a type of coastal bathymetric feature that have not been studied as extensively as other coastal features such as seamounts, headlands, and capes. Shoals are naturally-generated sandbars and ridges that are submerged in waters located along the coast. Shoals are found on top of continental shelves, and the morphology of shoals could be a potential hazard in maritime transport. Shoals are formed by the gradual accumulation of sand and convergent sediment transport, which can be influenced by both tidal currents and waves \cite{ashton01}. Shoals are susceptible to morphologic change over time, which could be accelerated by changes in sea level, which can further drive changes in coastal circulation and sediment transport. 

For this study, we investigate the extent to which shoal bathymetry impacts coastal ocean dynamics as it interacts with flow. Specifically, the impact of shoal bathymetry on offshore transport and mixing can provide new insight to near-coastal circulation dynamics and the role of shoals in mediating the exchange of water between the coast and the open ocean. 

The Diamond Shoals extend 13 kilometers out from Cape Hatteras, North Carolina, with a cluster of shallow sandbars that vary in form and depth over time. The Diamond Shoals are infamously known as the "Graveyard of the Atlantic," where more than 600 ships have been sunk in this region. The Cape Hatteras coastal region, where the Diamond Shoals lie, are a site of convergent flow between cool, fresh waters from the Middle Atlantic Bight (MAB) north of the shoal, and warm, salty water from the South Atlantic Bight (SAB) south of the Diamond Shoals \cite{seim22}. Furthermore, the Gulf Stream transitions from a boundary-trapped current to a free jet at The Point, where Cape Hatteras shifts in orientation, forming a dynamic biogeographic boundary with implications on regional carbon transport and biological productivity \cite{lee91}. Mean circulation and variability taken from observational work on the Cape Hatteras coastal region reveal that the coastal region is influenced by both local wind forcing and variation in the Gulf Stream spatially and temporally \cite{han22}

Observational data of waves and currents off Cape Hatteras show that mean subtidal flows are parallel to the coastline on either side of the Cape Hatteras point \cite{kumar13}. Flows in Diamond Shoals follows along the submerged shoal, suggesting that the shoal serves as an extension of the coastline. Along with the Diamond Shoals, the coasts of North Carolina are home to two other shoals, Frying Pan Shoals located off of Cape Fear, and Cape Lookout Shoals located off of Cape Lookout. The presence of these shoals along different parts of the coastline, each with distinct bathymetry and flow pattern, underline the motivation for this study, and we use the Diamond Shoals as a reference case for investigating how shoal bathymetry impacts offshore transport and mixing. 

In this paper, we organize as follows. Section 2 introduces background observational data and model setup. Section 3 presents our parameter sweep and detailed analysis on offshore flux, eddy kinetic energy, and baroclinicity. In Section 4, we comment on overall shoal-fluid dynamics and present our results in broader context. 


% \section{A Very Long Section Title that Should be Split Across Two Lines in the Table of Contents}


% Duis autem vel eum iriure dolor in hendrerit in vulputate velit esse molestie consequat, vel illum dolore eu feugiat nulla facilisis at vero eros et accumsan et iusto odio dignissim qui blandit praesent luptatum zzril delenit augue duis dolore te feugait nulla facilisi. Lorem ipsum dolor sit amet, consectetuer adipiscing elit, sed diam nonummy nibh euismod tincidunt ut laoreet dolore magna aliquam erat volutpat.   

% Ut wisi enim ad minim veniam, quis nostrud exerci tation ullamcorper suscipit lobortis nisl ut aliquip ex ea commodo consequat. Duis autem vel eum iriure dolor in hendrerit in vulputate velit esse molestie consequat, vel illum dolore eu feugiat nulla facilisis at vero eros et accumsan et iusto odio dignissim qui blandit praesent luptatum zzril delenit augue duis dolore te feugait nulla facilisi.   

% Nam liber tempor cum soluta nobis eleifend option congue nihil imperdiet doming id quod mazim placerat facer